%!TEX root = ../main.tex
% Findings v6 (2026-10-18). Converted from findings_v6_2026-10-18_claude.md. No design table; Figure 2 holds two board details.
% Participant and sticky quotations are italicized. Bracket citations give [document, timestamp].
% Bold bracket notes marked VERIFY must be checked against records or recordings before the dependent claim is kept.

% IDC '27 findings, v6 (2026-10-18)
%
% Rewritten on the v6 outline and the widening-circle arc (2026-10-17). Quotes are from the ATLAS.ti
% transcripts and sticky documents, cited as [document, timestamp]. Repeated fillers ("like," "um,"
% "uh") are removed from quotations where they carried no meaning, and an ellipsis marks any other
% omission; this convention belongs in the Methods, not here. Names are pseudonyms; where a
% transcript could not separate speakers, the quote is attributed to the team. [VERIFY] marks a fact
% or attribution that must be checked against records or recordings before the claim that depends on
% it is kept. No design table; Figure 2 holds two board details that the prose reads.
%

\section{Findings}
\label{sec:findings}

Section~\ref{sec:sense-making} describes how teens made sense of the listening presentation (RQ1). Across their proposals (RQ2), we identified three kinds of spatial decision: what to foreground, which listening perspective to capture or reproduce, and how to distribute an experience among listeners. In the following sections, we describe each kind of decision, the purposes teens gave for it, and the concerns they expressed about it. This order is our exposition and does not describe a sequence the teens followed. It moves from a decision about one listener's attention to decisions that reach other people, and the question of who makes the decision becomes harder to answer at each step. Throughout, we distinguish what teens reported from prior experience, how they responded to the presentation, and what they imagined in the design activity, and we note where an adult first proposed an idea. We do not report counts of participants, because four of the nine transcripts did not separate speakers reliably (Section~\ref{sec:methods}). Where a speaker could not be identified, we attribute the quote to the team.

\subsection{Hearing a Difference and Deciding Whether It Matters}
\label{sec:sense-making}

Teens' responses to the presentation involved three judgments that did not move together: whether they noticed a difference between the two versions, whether the difference changed their appreciation of the song, and whether they wanted anything about their listening changed. The spatial version carried a binaural mix, a moving display of circles marking where each sound was placed, the lyrics, and the video reduced to one side of the screen, and the two versions were not verified to be matched in loudness \textcolor{red}{\textbf{[VERIFY]}}. Middle schoolers listened in person on headphones \textcolor{red}{\textbf{[VERIFY: lab-provided ATH-M45 or participants' own]}} and high schoolers at home over Zoom on equipment we did not verify.

High schoolers in all three teams described the difference as small enough that they had to work to find it. Six of the eight first-question stickies in HS-H1 opened with a hedge of the form \textit{``I didn't hear too much of a difference''} \textcolor{red}{\textbf{[HDS1 Artifacts; VERIFY that the eight stickies are eight participants]}}. Amara (all names are pseudonyms), in HS-H1, reflected:

\begin{quote}
\textit{I was personally expecting more of a difference between the normal one and the spatial. There were some parts\ldots{} that I could hear only on my left ear and not in my right, and that was kind of cool. But you had to listen pretty carefully. And I actually went back to the first one to really see the differences, because initially, I was like, ``These literally sound exactly the same.''} [HDS.1, 37:32]
\end{quote}

Nora, in HS-H3, heard \textit{``a little bit\ldots{} I didn't really hear that much of an effect''} [HDS.3, 24:19], and Dev, in the same team, heard the spatial version as \textit{``mainly just louder and a bit more clear''} [HDS.3, 01:48] and preferred the standard one [16:05]. The middle schoolers hedged less. Their stickies were also shorter, handwritten, and prompted by a facilitator, so we do not read the contrast as a difference in what the two cohorts heard. One middle schooler in MS-L1 explained her preference for the spatial version: \textit{``I could hear this was coming from here, and this was coming from here. And without it\ldots{} everything was coming from one spot''} [MDS1.1, 07:04].

Across both cohorts, teens described the moving display as part of how they heard the spatial version. Sam, a musician in MS-L3, closed his eyes and found that \textit{``without really being able to see the dots show up\ldots{} it was kind of hard to differentiate what was on the left and what was on the right''} [MDS1.3, 08:59--09:15]. Mara, in HS-H3, wondered whether the difference she heard was \textit{``just a placebo effect from watching the video and the icons of the different music parts showing up''} [HDS.3, 08:20]. The facilitator replied, \textit{``I think that's not placebo, though. I think that's actually the spatial audio working''} [HDS.3, 08:41]. Mara then offered a second reading, that \textit{``seeing the audio mapped out on the screen makes you more aware of it as you hear it''} [HDS.3, 08:49], and added that \textit{``even without looking at the visual, I could definitely hear a difference\ldots{} it had more depth to it''} [HDS.3, 09:02]. It is difficult to know from our data whether the display made the mix audible or whether teens heard the mix and used the display to name what they heard. Mara's second reading came after the facilitator's reassurance, and we report it as hers.

For the teens who heard a difference, what changed was where their listening went. In both cohorts, teens described hearing the parts of the arrangement separately where before they had heard \textit{``one straight thing of audio''} [MDS1.3, 15:05]. A sticky in MS-L2 read \textit{``you could tell who was singing,''} and Mara heard \textit{``each instrument\ldots{} instead of it all blending together into one''} [HDS.3, 09:30]. Hearing the parts did not consistently change what teens thought of the song. HS-H1's own grouping of its stickies put seven under \textit{``Changed the Experience''} and two under \textit{``No Change in the Experience,''} and several of the seven were hedged. Nadia, in HS-H1, reflected: \textit{``at the end of the day, it's still the same music with the same notes\ldots{} I'd still listen to either version''} [HDS.1, 45:16]. Familiarity with the song cut both ways. A middle schooler in MS-L2 \textit{``saw the difference, but it just didn't change my overall feel,''} because \textit{``it's not a song I'm familiar with''} [MDS1.2, 04:32--04:52], whereas Amara, \textit{``someone who doesn't listen to K-pop,''} found that the spatial version \textit{``made me appreciate it more 'cause I was able to hear some of the unique sounds in it''} [HDS.1, 46:14].

Appreciating the difference did not mean wanting anything changed. In MS-L1, one participant's sticky read \textit{``Wouldn't change anything about my music or watch,''} and in discussion she repeated it: \textit{``I think I like how it is''} [MDS1.1, 13:55--14:04]. Another in the same team, asked what feature would make listening on the bus or during homework more fun, answered, \textit{``I think it's already fun''} [MDS1.1, 12:16]. A middle schooler in MS-L2, asked whether he would want spatial audio in games, said \textit{``Absolutely,''} then added, \textit{``Although I do have a headset that automatically does that\ldots{} I'm already doing this''} [MDS1.2, 16:01--16:09]. Teens judged the presentation against formats and places they already knew. Lena, in HS-H3, whose headphones were dead, judged the presentation \textit{``from what I've noticed from past audios that were like that,''} meaning 8D audio, which she found \textit{``sometimes\ldots{} a little bit too much''} [HDS.3, 21:56]. Concerts were the reference that recurred most. Sam explained that there \textit{``you kind of get to feel it all around you''} [MDS1.3, 14:54--15:05], and two high schoolers in HS-H3 found the presentation short of \textit{``the ground shaking''} \textcolor{red}{\textbf{[HDS.3, 28:59--29:56; VERIFY speaker]}} and of \textit{``a social atmosphere''} [HDS.3, 35:17]. Max, in MS-D3, applied the same kind of test to his own team's ideas: \textit{``what makes this sticker speaker special from just a small speaker?\ldots{} So I feel like the sticker thing's a cool idea\ldots{} but not the most''} [MDS2.3, 12:04--12:40]. These quotes illustrate how teens treated a spatial effect as something that had to earn its place against listening they already found adequate.

Responses to the demonstration did not line up with whether teens wanted to change their listening. In their proposals, they considered spatial possibilities in relation to particular activities, and the one change the demonstration reliably produced, hearing the parts of an arrangement separately, is where those proposals began.

\subsection{Separating Sounds for Attention and Presence}
\label{sec:foreground}

Teens proposed separating sounds in space so that one voice or sound could be followed among others. In one developed case, the separation proposed for following a conversation also changed how a teen imagined another person's presence. Nadia, in HS-H1, began from a problem with group calls and reasoned through a proposal in a single turn:

\begin{quote}
\textit{Whenever me and my friends are on FaceTime, it's like everyone's talking over each other and stuff like that. And I thought a cool way for spatial audio to help that is if they could make it so that it sounds like they're coming from different areas of the room. So if I'm sitting at my desk, one of my friends is in front of me to the left, and one of them is in front of me to the right, so that it feels like they're in\ldots{} I guess it'd kind of be weird to feel that they're in my room with me, but it would make it feel like that and easier communication.} [HDS.1, 55:21--55:41]
\end{quote}

Here, we see that the placement proposed for intelligibility and the presence Nadia found \textit{``weird''} are two consequences of one arrangement. Placing each friend's voice at a point in the room would make the call easier to follow. The same placement would make absent friends present in a room they are not in. Nadia noticed this mid-sentence, named it, and returned to the benefit. This sequence reads to us as ambivalence about the arrangement, since she did not withdraw the proposal.

Other proposals to separate sounds tied direction to a task the teen already had, and none of them carried the presence that Nadia noticed. Nadia herself recalled a museum audio guide where \textit{``the room is so big''} that she did not know \textit{``what it's talking about,''} and wanted \textit{``some sort of sense of direction of where it's talking about''} [HDS.1, 56:16]. In HS-H3, Dev imagined earbuds that follow the gaze: \textit{``you look at someone, then their voice would\ldots{} focus only on their voice, and if you look away, then it would kind of weaken''} [HDS.3, 52:39]. Lena connected that idea to her grandmother's hearing aids, where \textit{``everything is kind of the equal amount of sound that she's getting just amplified''} [HDS.3, 53:24--54:04]. In the middle school wrap-up, Sam told the room that in \textit{``a shooter game, if you hear someone behind you, you're gonna turn around''} [MDS1.3 wrap-up, 33:31]. In each of these proposals the separated sound belongs to a guide, a game, or a person the listener is already facing, so bringing it forward brings nobody new into the room. The hearing-aid benefit is an imagined one, and nothing in this section was built or tested.

The proposals differ in who would place the sounds. Dev's earbuds leave the decision with the wearer at every glance. Nadia's call does not say who would seat her friends around her desk, and none of the proposals considered the people being placed.

\subsection{Choosing a Listening Perspective to Capture or Reproduce}
\label{sec:perspective}

Two middle school design teams treated the perspective from which sound is captured or reproduced as something to choose, and the two cases differ in what the choice preserves. Both teams worked under the music-focused middle school prompt, and both designs took shape after an adult proposal, so we mark the adult contributions as they occur.

MS-D3's hoodie began when the team chose a sweatshirt from its own list and an adult suggested combining it with a teammate's \textit{``spatial recorder''} idea [MDS2.3, 12:57]. One teammate reframed the hoodie around a singer who \textit{``had a good song idea, they could press something, and it would record\ldots{} and you can play it again with spatial audio\ldots{} so you don't forget what you were saying''} \textcolor{red}{\textbf{[MDS2.3, 13:19--13:37; VERIFY speaker, merged label]}}. Others gave its purpose as a concert \textit{``you don't want to always have a thing to relive''} \textcolor{red}{\textbf{[MDS2.3, 13:51; VERIFY wording against the recording]}}. A teammate placed the mics \textit{``on each ear so you can capture the same stuff in here''} [MDS2.3, 19:07], and Max added a second mic near the mouth so that the wearer could \textit{``choose which you like better. Do you like it from right out of your mouth, or do you like it how you hear it better?''} [MDS2.3, 21:25] (Figure~\ref{fig:details}a). When an adult asked whether a friend could share the experience by putting the hoodie on, a teen answered, \textit{``This was how I experienced the concert''} [MDS2.3, 16:15--16:19]. Evident in Max's question is a distinction between two things a recording can keep. A mic at the mouth keeps the song idea, and mics at the ears keep how the wearer heard the concert. Only the second requires spatial capture. The team imagined replaying and sharing both.

MS-D2's practice app began with the facilitator's proposal to make an acoustic instrument spatial and his question, \textit{``are y'all okay if we start running with this idea?''} [MDS2.2, 06:16--07:44]. Within that frame, Rosa proposed \textit{``two settings''} [MDS2.2, 16:51], one in which \textit{``you could hear it run around or something''} [17:21] and one in which \textit{``you would sound like at a concert''} [18:12]. The facilitator then proposed hearing yourself as a Carnegie Hall audience would and introduced room acoustics [18:17--20:18]. Eli reduced the idea to \textit{``Big room, small room''} [20:13] and proposed that the app \textit{``generate a 3D room. And then maybe you could place yourself anywhere in that room''} [21:13--21:19], and Rosa added that a player could \textit{``search up a place''} [21:57--22:01]. The listener-position choice in this design is a playback choice for practice. Its framing as hearing oneself as an audience would is the facilitator's. The teens' additions were the two settings, the room sizes, and the room search.

The two cases differ in what the choice preserves. The hoodie keeps how an event was heard by the person who was there, and the decision about what to keep and whom to share it with sits with the wearer. The practice app places a performer in a simulated room, and the decision about which room sits with the performer, inside a design an adult framed. Both are middle school designs made under a prompt about music, and we do not claim the distinction for the high school teams.

\subsection{Distributing an Experience Among Listeners}
\label{sec:distribute}

Four designs distributed a spatial experience among listeners in different ways, and teens' reasons were about particular people in particular seats. HS-H2's \textit{``Virtual Hangout Space''} began from the teens' own FaceTime idea in the listening half of the session. Ren observed that \textit{``people my age video call so much. And if different people came from different sides, it would help with feeling like you're actually having a conversation''} [HDS.2, 28:54]. The design let \textit{``people choose what spot to sit at the table in the FaceTime app, and then their voices will project from there''} [HDS.2, 50:45]. The premium features that followed came from teens and facilitator in turn, and we attribute them turn by turn. Ben proposed a \textit{``Surround Sound wrapped''} [HDS.2, 01:01:32]. Hana proposed paying \textit{``a premium to claim that seat so nobody else can take it''} [01:03:03]. The facilitator proposed making \textit{``the table longer and longer''} for a friend who \textit{``said something you didn't like''} and a \textit{``silent little cone''} [01:03:40, 01:04:18], and Hana added, \textit{``Or a timeout?''} [01:04:30], recorded on the board as \textit{``Time out of shame.''} A further line on the board, \textit{``Muting people or changing their volumes,''} has no author in the transcript \textcolor{red}{\textbf{[VERIFY]}}. Hana called the set \textit{``really fun social gimmicks''} [HDS.2, 01:10:44], and the phrase was approving. Choosing a seat at this table changes where everyone else hears one's voice, and claiming a seat removes that choice from someone else. This difference did not come up in the team's discussion, and the board does not say how a claim would work.

HS-H3's car system was drawn by the facilitator (the first author) to the teens' instructions, and the questions about where the effect should go were his. When he asked whether the system should use headphones, Lena answered, \textit{``Headphones in the car?''}, and a teammate called the idea \textit{``insane''} and \textit{``unsafe''} \textcolor{red}{\textbf{[HDS.3, 01:07:13--01:07:37; VERIFY speaker]}}. When he asked whether she wanted her seat to shake, Lena said, \textit{``that sounds cool, but I don't know if I want to be driving on the interstate and all of a sudden the song comes on and I'm shaking''} [01:11:13--01:11:29], and a teammate moved the woofers to \textit{``the passengers only''} [01:12:08]. Asked later whether passengers might have headphones, the same teammate allowed that \textit{``keeping headphones optional would be nice''} [01:16:55--01:17:18]. Lena also wanted \textit{``two speakers in the middle so that the driver is having the same experience as the people in the rest of the car''} [01:13:17--01:13:38], which the board labels \textit{``A middle speaker to standardize experiences''} (Figure~\ref{fig:details}b). These were answers to facilitator questions. They do show the team placing the effect by role, with the shaking for the passengers and the same music for everyone.

The middle schoolers in MS-D1 designed for listening alone, and their reasons were also about the other people in the room. Owen worked through whether everyone in his room of speakers would hear the same thing, decided that \textit{``everyone probably wouldn't experience the same amount of sound,''} and proposed \textit{``doors for every person. Then there's speakers inside, and they have their own TV\ldots{} Alone. Well, soundproof walls''} [MDS2.1, 13:18--13:56]. Ivy's sheet-music tool played back \textit{``into speakers on your ears, so then you can hear it. But then others can't hear it. So it doesn't disturb anybody''} [MDS2.1, 22:03--22:38]. Owen and Ivy wanted to keep their sound to themselves, whereas Lena wanted everyone in the car to get the same sound. Payment entered the high school discussions as a condition on who gets the feature at all. Kat, in HS-H1, argued that spatial audio should come with the product, \textit{``'cause when you buy AirPods, noise cancellation is a feature in it''} [HDS.1, 01:36:29--01:37:20], then allowed \textit{``a premium 'cause it's a new thing''} with \textit{``a one-month free trial''} [01:37:38].

Across the four designs, the teens placed the decision about who hears what differently. The table's premium list places seat-claiming with whoever pays. The car team decided for the driver, and the facilitator's questions set that decision up. Owen's doors give each person their own sound, and Lena's middle speaker gives everyone the same. The teens did not say, for any of these arrangements, how one listener's choice would be weighed against the choices of the others who shared it, and we return to that question in Section~\ref{sec:discussion}.

\begin{figure}[t]
  \centering
  % Replace the placeholders with crops from the original boards:
  % \includegraphics[width=0.48\columnwidth]{hoodie-mics} \hfill \includegraphics[width=0.48\columnwidth]{car-speakers}
  \fbox{\parbox[c][0.18\textheight][c]{0.46\columnwidth}{\centering (a) Placeholder: MS-D3 hoodie, mic positions}}\hfill
  \fbox{\parbox[c][0.18\textheight][c]{0.46\columnwidth}{\centering (b) Placeholder: HS-H3 car, speakers}}
  \caption{Details from two design boards. (a) MS-D3's hoodie, with a mic at each ear and a second mic near the mouth, the two capture positions Max asked the wearer to choose between. (b) HS-H3's car, with woofers for the passengers and a middle speaker \textit{``to standardize experiences.''} \textcolor{red}{\textbf{[VERIFY: crops to be made from the original boards; check that the drawn mic positions match the transcript.]}}}
  \label{fig:details}
\end{figure}
